\documentclass{article}

\usepackage[utf8]{inputenc}
\usepackage[onehalfspacing]{setspace}
\usepackage{ngerman} % hiermit werden deutsche Bezeichnungen genutzt und 
  
\title{Exposé: Configuration-based mobile applications}
\author{Samuel Schepp}

\date{\today}

\begin{document}
	
\maketitle
	
\section*{Context}

Fabrik19, B2B mobile and terminal (stele) apps.

configuration based apps.

how? download layout, data definition, actions, metadata.

all apps use the same container, which is always up to date and uses latest technology.

some metadata has to be baked in on compile time.

\section*{Problem definition}

why configuration based?

long waiting times on app store approval

example: customer terminal, city marketing stele: lot of work to redeploy large amount of units.

lot of different apps need same boilerplate code. hard to migrate all to newest technology.

\section*{Motivation}

use of native mobile features, camera, background gps, beacons, push notification (with payload). caching, fast ui, good battery life.

no need to depend on third party frameworks, direct use of phone vendor's apis.

why no code generator?

no benefit, impossible development/debug process, unsuitable for highly dynamic applications.

\section*{Distinction}



\end{document}